% Section 12.3, from lectures/L12/README.md: what you should be able to do afterwards, the
% questions to test yourself, and the next lecture.

\section{Review}
\label{c12:sec:review}

\needspace{6\baselineskip}
\subsection*{What you should be able to do now}
After this chapter, you should be able to:
\begin{itemize}
  \item Define a real-time system without using the word ``fast'', and say why an average is not
    evidence.
  \item Name three sources of latency and say, for a given piece of code, which one it contributes
    to.
  \item Explain priority inversion, and say what priority inheritance does about it.
  \item Choose a preemption model for a stated requirement and defend it.
  \item Say what \code{PREEMPT_RT} does to a \code{spinlock_t} and to an interrupt handler.
  \item Say why \code{raw_spinlock_t} exists and when using one is a mistake.
  \item Find, in unfamiliar driver code, the four habits that make it hostile to real time.
  \item Run \code{cyclictest}, read its output, and state precisely what it did and did not measure.
  \item Present a latency measurement together with the conditions that make it meaningful.
\end{itemize}

\needspace{6\baselineskip}
\subsection*{Questions to test yourself}
\begin{enumerate}
  \item System A responds in 50 us on average and 30 ms at worst. System B responds in 5 ms every
    time. Which is real time, and for what requirement?
  \item Your handler takes a \code{spinlock_t} and does 200 us of work. What changes under
    \code{PREEMPT_RT}, and is the change an improvement?
  \item Why can \code{PREEMPT_RT} not simply make every lock a sleeping lock?
  \item \code{cyclictest} reports a maximum of 40 us. What does that tell you about your device's
    interrupt latency?
  \item A colleague reports ``we measured 12 us latency''. What four questions do you ask before
    believing it applies to your product?
  \item Why is a \code{udelay(50)} in a driver a real-time problem even though 50 us is short?
  \item You measure lower latency under QEMU than the datasheet claims for real silicon. What
    happened?
\end{enumerate}

\needspace{6\baselineskip}
\subsection*{Where to look}
\begin{itemize}
  \item \file{tools/qa-dev/qa-dev.c} shows exactly when the timestamp is taken, which is the one
    thing you must be sure of before trusting a latency measurement built on it.
\end{itemize}

\needspace{6\baselineskip}
\subsection*{Looking ahead}
This is the last chapter, and with it the book's main text is complete. It ends with a driver that
probes off a device tree, services interrupts, blocks a reader correctly, registers with two
subsystems, and whose latency you have measured and can describe honestly. What remains is to work
the exercises in \secref{c12:sec:exercises} and check the ones that are not code against the
solutions in Appendix~\ref{sol:ch:solutions}. After that, the two written papers in
Appendices~\ref{exam:ch:about} to \ref{ansb:ch:answers} are there to test, on paper and with nothing
in front of you, how much of all twelve chapters you can reconstruct on your own.

What to do next, in rough order of usefulness: read a real driver in \file{drivers/} for a device
you own and see how much of it you now recognise; write a driver for hardware that exists; and send
a patch, because the review you get on a first patch to a subsystem mailing list is worth more than
any course.
